\section{Access Path Trees}\label{sec:apt}

An access path describes how the program reached a particular
substructure of Move's state: it starts at some named entry
point---a local of the current frame, a globally-stored resource, or
a reference parameter passed in by the caller---and follows a
sequence of field projections, variant projections, and indexings.
Every borrow in the program produces a reference that can be
described by its access path. The runtime records these paths
explicitly: it groups them into a forest of trees, one per entry
point, and attaches each live reference to the node at the end of
its path.

A reference is \emph{valid} when the path it represents still leads
to a reachable, type-compatible substructure of the current state.
The runtime approximates this by marking references invalid when
destructive events happen along their paths; the rules in
Sect.~\ref{sec:semantics} spell out which events invalidate which
references.

\paragraph{Derivation.}
A reference $\refn{2}$ is \emph{derived from} $\refn{1}$ when its
access path is built by extending the path of $\refn{1}$ through one
of the language's borrowing operations: field projection, variant
projection, indexing, freezing, or being passed across a function
call as a reference parameter. In the access path tree this comes
out simply: $\refn{2}$ is derived from $\refn{1}$ exactly when its
node lies in the subtree below $\refn{1}$'s. The full derivation
order is the reflexive-transitive closure.

\paragraph{Roots and edges.}
Each tree's root is one of the named entry points mentioned above.
There are three kinds:
\begin{itemize}
\item $\local{i}$, the $i$-th local of the current frame;
\item $\glob{\tau}$, the global resource of type $\tau$ (we abstract
addresses);
\item $\param{i}$, the abstract value behind the $i$-th reference
parameter passed in by the caller.
\end{itemize}
Edges carry labels: a field offset for record fields, a variant
field for variant projections, and a single distinguished label $0$
for vector elements. The vector label is an over-approximation: rather
than track each index, we collapse all of them. This costs precision
on element-wise borrows but keeps the trees small.

Trees are built lazily. The runtime never materializes a node until
the program actually borrows through it.

\paragraph{State.}
A runtime state is a tuple $\sigma = (T, R, S, F)$.
\begin{itemize}
\item $T : \mathrm{Root} \rightharpoonup \mathrm{APT}$ maps each root
to its tree. The map is partial: roots that have not been borrowed
have no tree.
\item $R$ assigns each live reference identifier a triple $(\mu, \pi,
n)$. Here $\mu \in \{\mut, \imm\}$ is the mutability, $\pi \in
\{\Live, \Pois\}$ is the validity flag (we call $\Pois$ \emph{poisoned}
for short), and $n = \node{\mathrm{root}}{\mathrm{id}}$ is the
reference's \emph{qualified node}---the tree node it is attached to.
\item $S$ is a shadow stack mirroring the machine's operand stack.
Each entry is either a non-reference token or a reference identifier.
\item $F$ is the per-frame stack. Each frame carries its own $T$ and
$R$, a vector of shadow locals, and a map from this frame's
reference-parameter indices back to qualified nodes in the caller,
which is what return transport uses.
\end{itemize}

\paragraph{Aliasing as a tree relation.}
Two references alias when one's node is an ancestor or descendant of
the other's. The runtime decides this by walking parent pointers, in
time linear in the depth of the tree. Each node also holds the set of
references attached to it, so an invalidation walk only has to visit
the affected subtree---it does not have to scan the whole reference
table.

\section{Invalidation Rules}\label{sec:semantics}

We give the rules as a small-step relation $\sigma \to \sigma'$
indexed by the bytecode instruction being executed. Rules for
operations that do not touch the shadow state---constants,
arithmetic, stack reshuffles---are omitted.

Every non-trivial rule has the same shape. An event destructively
changes the substructure reachable through some node $n$, and the
runtime walks the tree from $n$ to poison every reference attached to
nodes related to $n$. Different events pick different sets of
related nodes (self, ancestors, descendants, or unions) and different
mutability filters. We write $\mathrm{invalidate}(N, \phi)$ for the
operation that poisons every reference attached to some node in
$N$, filtered by $\phi \in \{\mut, \imm, \mathsf{all}\}$. A poisoned
reference cannot be borrowed from, read, written, frozen, or passed
as a call argument; any such attempt is a runtime safety violation.

\paragraph{Borrow, read, freeze.}
Borrowing a local $i$ creates a fresh reference attached to
$\local{i}$, materializing $T(\local{i})$ if it does not yet exist. A
field borrow with label $\ell$ from a reference $\refn{}$ at node $n$
consumes $\refn{}$ and creates a fresh reference at the child
$n.\ell$. \texttt{ReadRef} pops the top reference after a validity
check; the node is unchanged. \texttt{FreezeRef} on a live $(\mut, n)$
reference consumes it and produces a fresh $(\imm, n)$ reference at
the same node---an immutable re-view of the same path. None of these
events invalidates anything.

\paragraph{Destructive write.}
\texttt{WriteRef} through a mutable reference at $n$ overwrites the
value at $n$, which affects both its ancestors (whose contents have
changed) and its descendants (which may have been replaced). The rule
is:
\[
\mathrm{invalidate}\bigl(\{n\} \cup \Desc(n) \cup \Anc(n),\ \imm\bigr),\qquad
\mathrm{invalidate}\bigl(\Desc(n),\ \mut\bigr).
\]
The first line poisons every immutable view of $n$, of any ancestor,
and of any descendant. The second poisons every mutable view of a
descendant. Mutable views of $n$ itself are deliberately spared: a
program that consumes one mutable view and immediately reborrows from
$n$ stays accepted, which matches static borrow checking.

\paragraph{Move, overwrite, variant mutation.}
Moving or storing over local $i$ replaces the value at $\local{i}$
outright, so we invalidate every reference reaching through it,
regardless of mutability:
$\mathrm{invalidate}(\{r_i\} \cup \Desc(r_i),\ \mathsf{all})$.
Variant mutation is the same rule applied at the discriminator:
references derived from the prior variant's projection sit in
descendant nodes, and the change to the discriminator poisons them.
The enum example from Sect.~\ref{sec:intro} falls out of this case.

\paragraph{Calls.}
A call $f(\overline{a})$ with reference arguments
$a_{i_1}, \ldots, a_{i_k}$, each at node $n_j$ with mutability
$\mu_j$, proceeds in four steps.
\begin{enumerate}
\item Check each $\refn{j}$ is live. Then lock the subtree at $n_j$:
exclusively if $\mu_j = \mut$, shared otherwise. An exclusive lock
conflicts with any other lock on a related node; a shared lock
conflicts only with an exclusive one. Any conflict is a violation.
This is what catches overlapping mutable arguments from
Sect.~\ref{sec:examples}.
\item For each mutable argument, apply the destructive-write rule at
$n_j$. The callee may write through the reference, so the discipline
treats the call itself as a write.
\item Allocate a new frame. The callee's reference parameters become
fresh references attached to parameter roots $\param{j}$. The new
frame records, for each $j$, the qualified node $n_j$ in the
caller---this is what return transport (below) uses.
\item Release the locks on the caller's subtrees.
\end{enumerate}
Fig.~\ref{fig:lock} illustrates a conflicting pair.
\begin{figure}[h]
\centering
\begin{minipage}{0.92\textwidth}\small
Calling \texttt{foo(r1, r2)} with $r_1 = \&\mut x$, $r_2 = \&\mut x.g$:
\begin{verbatim}
  root                       r1 takes exclusive lock on subtree(x).
   +-- x     <-- locked      r2 takes exclusive lock on subtree(x.g),
        +-- f                which is a descendant of x: lock conflict
        +-- g  <-- locked    ==> runtime safety violation.
\end{verbatim}
\end{minipage}
\caption{Call-boundary subtree locking.}\label{fig:lock}
\end{figure}

\paragraph{Returns.}
On return, every reference on the shadow stack must trace back to a
reference parameter of the current frame. Walking the parent pointers
of its qualified node should land on some $\param{j}$. A reference
rooted at a local or a global would dangle past the frame's lifetime
and is rejected. Each surviving reference is then transported into
the caller: starting from the caller's qualified node for parameter
$j$, the runtime walks the same labels and attaches a fresh reference
at the resulting node. Returning two mutable references that share a
subtree is caught by the same exclusive-lock pass used at calls.

\paragraph{Opaque effects.}
Some operations have bodies the runtime cannot see: native functions
implemented outside the bytecode language, closure captures, and
dynamic-dispatch targets. We handle them through a small interface:
each opaque operation declares, per returned reference, which input
reference argument it derives from. At a call site, the runtime
conjures a fresh child of that argument's node and attaches the
returned reference to it; ordinary call-boundary checks still apply.
The interface is part of the trusted base---an incorrect declaration
silently weakens the guarantees. We return to this in
Sect.~\ref{sec:discussion}.
